% ATTEMPT_STATUS: unresolved
% ATTEMPT_ORIGIN: new research, 2026-10-01
% TOP_PROBLEM: {"id": 30006806, "title": "Universal Optimality of the Hexagonal Lattice in Dimension Two", "category_id": 6, "category": {"id": 6, "name": "geometry", "display_name": "Geometry", "description": "Euclidean and non-Euclidean geometry, geometric structures.", "slug": "geometry", "order_index": 6, "created_at": "2026-07-31T15:26:25.671Z"}, "status": "open"}
% FIRST_POSTED: 2026-10-01
% !TeX program = lualatex
% Definition and sources only; this file is not an LLM solution attempt.
\documentclass[11pt]{article}
\usepackage[margin=25mm]{geometry}
\usepackage{fontspec}
\setmainfont{Latin Modern Roman}
\usepackage{amsmath,amssymb,mathtools}
\usepackage{unicode-math}
\setmathfont{Latin Modern Math}
\usepackage{xurl,hyperref}
\hypersetup{colorlinks=true,urlcolor=blue}
\setcounter{secnumdepth}{0}
\setlength{\parindent}{0pt}
\setlength{\parskip}{5pt}
\setlength{\emergencystretch}{3em}
\begin{document}
\section*{\texorpdfstring{Universal Optimality of the Hexagonal Lattice in Dimension Two}{Universal Optimality of the Hexagonal Lattice in Dimension Two}}\label{30006806}
\subsection{Definitions and mathematical statement}
Let $A_2$ be the triangular lattice rescaled to covolume one. A competitor $C\subset{\mathbb{R}}^2$ is discrete and closed with asymptotic density one in centered balls. A function $f:(0,\infty)\to{\mathbb{R}}$ is completely monotone when it is smooth and $(-1)^kf^{(k)}\ge0$ for every $k\ge0$. With $p(r)=f(r^2)$, define the lower energy by
\[
 \underline E_p(C)=\liminf_{R\to\infty}
 \frac{\sum_{x,y\in C\cap B_R,\ x\ne y}p(|x-y|)}{|C\cap B_R|}.
\]
The assertion is $\underline E_p(C)\ge E_p(A_2)$ for all these $C,f$, allowing infinite energies. The sum counts ordered pairs, and the lattice energy uses the same convention. Competitors need not be lattices or periodic.
\par\smallskip\textit{Formulation and concept references:} \href{https://arxiv.org/abs/1902.05438}{[S1]}.

\subsection{Short English statement}
Does the density-one triangular lattice minimize every completely monotone squared-distance interaction energy among all planar point configurations of the same asymptotic density?
\subsection{Research attempt}
\textbf{Status: UNRESOLVED.} We justify passage from Gaussian lower bounds to positive mixtures, compute the Gaussian lattice-energy formulation, and prove that the normalized triangular lattice beats the square lattice for sufficiently short-range and sufficiently long-range Gaussians. These comparisons cover neither all Gaussian scales nor all density-one competitors.

\paragraph{Source and dimensional scope.}
Cohn--Kumar--Miller--Radchenko--Viazovska [S1], primary abstract checked 1 October 2026, proves universal optimality in dimensions eight and twenty-four. It does not prove the two-dimensional target. We use its energy framework as context and analyze simple two-dimensional competitors directly. All sums here use ordered pairs, matching the definition; changing to unordered pairs would divide both compared energies by two.

\paragraph{Gaussian energies of a lattice.}
For $a>0$ let $p_a(r)=e^{-\pi ar^2}$ and, for a covolume-one lattice $L$, set
\[
E_a(L)=\sum_{v\in L\setminus\{0\}}e^{-\pi a|v|^2}.
\]
This sum converges: the number of lattice points in a disk of radius $R$ is $O_L((R+1)^2)$, by surrounding them with disjoint small disks, while the Gaussian decays faster than that growth. The finite-window ordered-pair energy converges to the displayed sum. Indeed for each fixed displacement $v$, the fraction of lattice points in $B_R$ whose translate by $v$ also lies in $B_R$ tends to one. Boundary points are confined to an annulus of fixed width and contribute $O_v(R)$ compared with $\pi R^2+O_L(R)$ points. The fraction is always at most one; dominated convergence over the summable displacement weights completes the argument.

\paragraph{A normalization check.}
Write the triangular lattice with basis $(s,0)$ and $(s/2,\sqrt3s/2)$, where $s^2=2/\sqrt3$. Its covolume is $s^2\sqrt3/2=1$, and
\[
|mu+nv|^2=s^2(m^2+mn+n^2).
\]
For integers $(m,n)\ne(0,0)$, the quadratic form is a positive integer, so its minimum is one and the shortest squared vector length is $\lambda_T^2=2/\sqrt3$. There are six such shortest vectors. The covolume-one square lattice has shortest squared length $\lambda_S^2=1$ and four shortest vectors. Comparing unscaled triangular and square lattices would confuse density with shape; the normalization above avoids that error.

\paragraph{Short-range Gaussian comparison.}
For every fixed lattice $L$ and $a\ge1$, convergence of its theta series gives a finite constant $C_L$ with
\[
E_a(L)\le C_L e^{-\pi a\lambda_L^2}.
\]
To prove this, factor out the shortest-vector exponential and bound the remaining sum by $\sum_{v\ne0}e^{-\pi(|v|^2-\lambda_L^2)}$, using $a\ge1$. On the other hand $E_a(\mathbb Z^2)\ge4e^{-\pi a}$. Consequently
\[
\frac{E_a(A_2)}{E_a(\mathbb Z^2)}
\le\frac{C_{A_2}}4 e^{-\pi a(2/\sqrt3-1)}\longrightarrow0.
\]
Thus the triangular lattice has strictly smaller energy than the square lattice for all sufficiently large $a$. This is a rigorous asymptotic comparison, although no optimized numerical transition value is computed.

\paragraph{Long-range comparison by Poisson summation.}
For a covolume-one lattice in dimension two, Gaussian Poisson summation gives
\[
1+E_a(L)=a^{-1}(1+E_{1/a}(L^*)).
\]
The square lattice is self-dual. Directly computing the dual basis of the normalized triangular basis shows that its dual has the same Gram matrix up to changing a basis vector's sign; it is an orthogonal image of $A_2$. Hence
\[
E_a(A_2)-E_a(\mathbb Z^2)
=a^{-1}\bigl(E_{1/a}(A_2)-E_{1/a}(\mathbb Z^2)\bigr).
\]
The large-parameter comparison therefore also proves strict superiority of $A_2$ for all sufficiently small positive $a$. Poisson summation applies because the Gaussian is Schwartz and the lattices have the stated covolume. It does not address nonperiodic competitors by itself.

\paragraph{What Gaussian inequalities would imply for mixtures.}
Suppose a nonnegative potential has a representation $p(r)=\int_0^\infty e^{-\pi ar^2}\,d\mu(a)$ by a positive measure, and suppose the desired Gaussian inequality has been established for a given competitor at every $a$. For finite windows, Tonelli's theorem interchanges the nonnegative sum and integral. Along any sequence of radii realizing the lower energy, Fatou's lemma gives
\[
\underline E_p(C)\ge\int_0^\infty\underline E_{p_a}(C)\,d\mu(a)
\ge\int_0^\infty E_a(A_2)\,d\mu(a)=E_p(A_2).
\]
The use of a liminf permits this inequality, but does not permit an unjustified equality exchanging limit and integral. For completely monotone functions, the general Laplace representation is a separate classical analytic theorem; this paragraph proves the energy passage for a supplied positive representation, including every finite positive Gaussian mixture. It does not reprove that representation theorem.

\paragraph{Remaining gap and verification.}
The square-lattice comparisons are proved symbolically, with no claimed numerical optimization. Positive constant potentials yield infinite energy for every density-one configuration because the number of other points per point tends to infinity; they do not distinguish the proposed optimizer. All nonnegative mixture manipulations remain valid with infinite energies.

The full target still needs Gaussian lower bounds at every scale against arbitrary discrete closed density-one configurations. Beating one lattice in two asymptotic regimes supplies neither statement. Even a theorem for every lattice would leave periodic sets with several points per cell and nonperiodic configurations untreated. The exact missing step is a sharp universal energy certificate with the canonical normalization and density hypotheses.

\subsection{Sources}
\begin{itemize}
\item[S1]Universal optimality of the E8 and Leech lattices and interpolation formulas,1902.05438v3,10June2022.
\url{https://arxiv.org/abs/1902.05438}.
\textit{Formulation source.}
\end{itemize}
\end{document}
